\input{estilo}
% El flujo de la pizarra: el frontend escribe en la base (con la tabla
% adentro); el worker hace receive contra la base y manda a SQS; la flota toma
% de SQS.
\begin{document}
\begin{figura}\begin{tikzpicture}
  \node[caja, minimum width=13mm, minimum height=12mm] (f) at (0,1.3) {FE};
  \pic[scale=1.4] at (0,-2.0) {base};
  \draw (-0.35,-0.85) rectangle (0.35,-0.6) (-0.35,-1.35) rectangle (0.35,-1.1) (-0.35,-1.6) -- (0.35,-1.6);
  \draw[msg] (f.south) -- (0,-0.15);
  \node[caja, minimum width=15mm, minimum height=17mm] (w) at (6.0,1.0) {worker};
  \draw[msg] (w.south) to[out=-90, in=0] node[pos=0.65, below=4pt] {receive} (0.85,-1.3);
  \node[cola, minimum height=22mm, minimum width=9mm] (q) at (9.0,1.0) {SQS};
  \draw[msg] (w.east) -- (q.west);
  \foreach \d in {0.3, 0.15} \node[caja, fill=white, minimum width=13mm, minimum height=13mm] at (12.0+\d,1.2-\d) {};
  \node[caja, minimum width=13mm, minimum height=13mm] (k) at (12.0,1.2) {};
  \draw[msg] (k.west) to[bend right=25] (q.east);
\end{tikzpicture}\end{figura}
\end{document}
